% Original analytic teaching diagram, 2026-11-12.
% Schematic (not to scale, KL is not a distance): target Pi outside the candidate set Q (yellow region);
% best candidate Q^dagger; iterates Q^(t) (green) approaching it; a larger set Q' (dashed) reaching closer to Pi.
% Two of the three error sources are drawn: expressiveness (Pi to Q^dagger) and optimisation gap
% (Q^(t) to Q^dagger); the third, error of the computed objective value, is discussed in the text.
\definecolor{VB}{HTML}{7188A3}\definecolor{VBF}{HTML}{E5EAF0}
\definecolor{VG}{HTML}{829C8B}\definecolor{VGF}{HTML}{E8EEE9}
\definecolor{VY}{HTML}{B4A16F}\definecolor{VYF}{HTML}{F1ECDD}
\definecolor{VR}{HTML}{A7787D}\definecolor{VRF}{HTML}{F0E5E6}
\definecolor{VP}{HTML}{9588A5}\definecolor{VPF}{HTML}{ECE8F0}
\definecolor{VI}{HTML}{3B4252}\definecolor{VM}{HTML}{6B7280}\definecolor{VL}{HTML}{4C4D4F}
\begin{tikzpicture}[x=1mm,y=1mm,>=Stealth,font=\figurefont\fontsize{8}{10}\selectfont,
lab/.style={text=VM,font=\figurefont\fontsize{7.5}{9}\selectfont}]
\draw[VL,line width=.5pt,dashed] (34,18) ellipse (40 and 17);
\fill[VYF] (24,16) ellipse (22 and 11);
\draw[VY,line width=.8pt] (24,16) ellipse (22 and 11);
\coordinate (P) at (64,26);
\coordinate (Qd) at (45.4,18.6);
\fill[VB] (P) circle (1.2);
\node[text=VB,anchor=west] at (65.5,26) {$\Pi$};
\draw[VR,line width=1.1pt,<->,shorten <=1.4mm,shorten >=1.4mm] (Qd) -- (P);
\node[text=VR,anchor=south] at (56,24.5) {表达限制};
\fill[VY] (Qd) circle (1.1);
\node[text=VI,anchor=north] at (45.4,17.2) {$Q^{\dagger}$};
\foreach \x/\y in {10/8,20/12,30/15,37/16.8,41.5/17.8} {\fill[VG] (\x,\y) circle (.8);}
\draw[VG,line width=.7pt] (10,8) -- (20,12) -- (30,15) -- (37,16.8) -- (41.5,17.8);
\node[text=VG,anchor=east] at (8.5,8) {$Q^{(0)}$};
\node[text=VG,anchor=south] at (30,16) {$Q^{(t)}$};
\draw[VR,line width=.8pt,<->] (31,21.5) -- (44.5,21.5);
\node[text=VR,anchor=south] at (37.8,22) {优化差距};
\node[lab] at (19,21) {候选集合 $\mathcal{Q}$};
\node[lab,anchor=south] at (34,35.5) {更大的集合 $\mathcal{Q}'$};
\end{tikzpicture}
